\documentclass[11pt]{article}

% ── Packages ──────────────────────────────────────────────────────────────────
\usepackage[margin=1in]{geometry}
\usepackage{microtype}
\usepackage{hyperref}
\usepackage{url}
\usepackage{booktabs}
\usepackage{tabularx}
\usepackage{listings}
\usepackage{xcolor}
\usepackage{amsmath}
\usepackage{amssymb}
\usepackage{graphicx}
\usepackage{caption}
\usepackage{subcaption}
\usepackage{enumitem}
\usepackage{titlesec}
\usepackage{parskip}
\usepackage{natbib}
\usepackage{authblk}
\usepackage{abstract}

% ── Listing style ─────────────────────────────────────────────────────────────
\definecolor{codebg}{RGB}{248,248,248}
\definecolor{codeframe}{RGB}{220,220,220}
\definecolor{kwcolor}{RGB}{0,0,180}
\definecolor{commentcolor}{RGB}{100,100,100}

\lstdefinelanguage{ICS}{
  keywords={ALLOW,DENY,REQUIRE,WITHIN,ON,WITH,UNLESS,IF,CLEAR},
  keywordstyle=\color{kwcolor}\bfseries,
  comment=[l]{\#},
  commentstyle=\color{commentcolor}\itshape,
  stringstyle=\color{black},
  basicstyle=\ttfamily\small,
  breaklines=true,
  frame=single,
  backgroundcolor=\color{codebg},
  rulecolor=\color{codeframe},
  showstringspaces=false,
  moredelim=[s][\color{kwcolor}]{"\#\#\#ICS:}{"\#\#\#},
}

\lstset{
  basicstyle=\ttfamily\small,
  breaklines=true,
  frame=single,
  backgroundcolor=\color{codebg},
  rulecolor=\color{codeframe},
  showstringspaces=false,
  numbers=none,
}

% ── Hyperref setup ────────────────────────────────────────────────────────────
\hypersetup{
  colorlinks=true,
  linkcolor=blue!60!black,
  citecolor=blue!60!black,
  urlcolor=blue!60!black,
}

% ── Title ─────────────────────────────────────────────────────────────────────
\title{\textbf{ICS: A Formal Instruction Contract Specification\\
for Large Language Model Interfaces}\\[0.4em]
\large Version 0.1 --- Technical Report}

\author[1]{[Author Name]}
\affil[1]{[Affiliation]}

\date{\today}

% ─────────────────────────────────────────────────────────────────────────────
\begin{document}
\maketitle

% ── Abstract ──────────────────────────────────────────────────────────────────
\begin{abstract}
Instructions to Large Language Models (LLMs) are typically authored as
informal prose, resulting in recurring failure modes: implicit constraints
that disappear across sessions, mixed-lifetime context that cannot be
cached, ambiguous output expectations that are resolved only by repeated
re-prompting, and prompt content that cannot be audited or validated
programmatically. We introduce the \emph{Instruction Contract Specification}
(ICS), a five-layer structured format that treats LLM instructions as formal
interfaces rather than natural language requests. ICS separates instruction
content by \emph{lifetime}: permanent context, capability declarations,
session state, per-invocation tasks, and output contracts. This separation
enables prompt-caching APIs to serve stable content at reduced rates. We
provide a formal grammar for capability directives, a conformance checker, and
empirical evaluation on a real-world B2B payments platform workload. Across
ten API invocations on \texttt{claude-haiku-4-5-20251001}, ICS-structured
instructions achieve a \textbf{77.8\%} reduction in effective API cost
compared to the na\"{i}ve flat-prompt baseline, with a proven lower bound of
55\% token savings at $N{=}10$ that is independent of model or counting
method. ICS is model-agnostic, framework-independent, and available as an
open specification.
\end{abstract}

% ─────────────────────────────────────────────────────────────────────────────
\section{Introduction}
\label{sec:intro}

The software industry has spent decades developing discipline around
\emph{interfaces}: API contracts, schemas, protocols, and type systems that
make the behavior of software components explicit, auditable, and
maintainable. These tools exist because interfaces that rely on shared
context, informal conventions, and implied constraints do not compose and do
not scale.

LLM instructions have not received the same treatment. Instructions are
typically written as natural language prose, structured loosely by convention,
and stored in files with no machine-readable semantics. When a system that
uses an LLM instruction misbehaves, diagnosing whether the fault lies in the
model, the instruction, or the caller's interpretation of the output requires
manual inspection. When the instruction must be reused in a new context or
maintained by someone who was not involved in writing it, implicit assumptions
collapse silently.

This produces four failure modes that recur across real deployments:

\begin{enumerate}[leftmargin=*]
\item \textbf{Context collapse.} Long-lived facts (system invariants,
architectural constraints) are mixed with short-lived context (session
decisions, per-task parameters) into a single prompt block. Callers cannot
determine what can be safely cached and what must be re-sent on every
invocation, resulting in either unnecessary token cost or stale-context bugs.

\item \textbf{Implicit constraints.} Team conventions---``never modify the
API layer,'' ``always return amounts in minor units,'' ``assume Python
3.11''---are documented in READMEs or exist only in the knowledge of the
original author. When the instruction is reused or transferred, these
constraints disappear silently.

\item \textbf{Ambiguous output expectations.} Instructions describe
\emph{what to do} without specifying what a correct result looks like.
Output evaluation becomes subjective, failure states are undetectable by
callers, and retry behavior is undefined.

\item \textbf{Retry as a first resort.} Ambiguous instructions produce
inconsistent outputs. The standard response is to rephrase and retry---which
treats ambiguity as a model failure rather than an interface failure, and
produces no durable fix.
\end{enumerate}

We argue that these failure modes share a root cause: \emph{LLM instructions
are not treated as interfaces.} The solution is not better prompting
intuition---it is applying the same engineering discipline to instructions
that is already applied to APIs, schemas, and protocols.

\paragraph{Contributions.} We present the \emph{Instruction Contract
Specification} (ICS), which makes the following contributions:

\begin{enumerate}[leftmargin=*]
\item A five-layer instruction format that separates content by lifetime
(\S\ref{sec:spec}), making stable context structurally distinct from
per-invocation content.

\item A formal directive grammar for capability declarations
(\S\ref{sec:capgram}), replacing prose constraints with machine-verifiable
directives.

\item A required output contract with enumerated fields, eliminating
ambiguous output expectations and undefined failure states (\S\ref{sec:oc}).

\item Empirical evidence that ICS's lifetime separation enables prompt-caching
APIs to reduce effective API cost by 77.8\% on a real-world workload, with a
proven token-saving lower bound that holds for any conformant instruction
(\S\ref{sec:eval}).

\item An open specification, reference validator, and worked examples across
five real-world domains (\S\ref{sec:examples}).
\end{enumerate}

% ─────────────────────────────────────────────────────────────────────────────
\section{Related Work}
\label{sec:related}

\paragraph{Prompt templates and structured prompting.}
Early systematic approaches to prompting introduced reusable templates and
labeling schemes~\citep{bach2022promptsource,sanh2022multitask}. DSPy
\citep{khattab2023dspy} compiles natural-language signatures into optimized
prompt programs, treating prompt construction as a differentiable pipeline.
LMQL~\citep{beurer2023lmql} and Guidance~\citep{guidance2023} add
programming-language constructs (loops, conditionals, type constraints) to
prompt authoring. TypeChat~\citep{typechat2023} uses TypeScript type
definitions to constrain LLM outputs to well-typed JSON. ICS differs from all
of these: it does not optimize prompts, does not compile them, and does not
constrain \emph{model decoding}. Instead, ICS defines a \emph{structural
contract} for the \emph{instruction itself}---what the caller commits to
providing and how it is organized---independent of the underlying model or
inference engine.

\paragraph{Chain-of-thought and instruction tuning.}
Chain-of-thought prompting~\citep{wei2022chain} and related
techniques~\citep{kojima2022zero,wang2022self} elicit multi-step reasoning
by structuring the prompt's \emph{content}. Instruction
tuning~\citep{ouyang2022training,wei2022finetuned} aligns model behavior with
natural-language instructions at training time. These methods address how
models \emph{reason}; ICS addresses how instructions are \emph{structured and
validated}. The two concerns are orthogonal: an ICS-compliant instruction can
include chain-of-thought examples in its TASK\_PAYLOAD layer.

\paragraph{Memory and context management.}
MemGPT~\citep{packer2023memgpt} introduces a hierarchical memory system that
pages context in and out of an LLM's context window at runtime, enabling
unbounded effective context. Retrieval-augmented generation
(RAG)~\citep{lewis2020rag} injects retrieved documents into the prompt at
query time. Both address \emph{what content} the model sees; ICS addresses
\emph{how that content is organized and classified by lifetime}. The two
approaches are complementary: a RAG system could use ICS to separate
retrieved content (TASK\_PAYLOAD) from stable system context
(IMMUTABLE\_CONTEXT).

\paragraph{Prompt caching and token efficiency.}
Anthropic's prompt caching API~\citep{anthropic2024caching} and analogous
mechanisms in other providers allow prefill KV-cache reuse for stable prompt
prefixes. Gim et al.~\citep{gim2023prompt} propose prompt cache architectures
for shared prefix reuse across requests. ICS's lifetime-separation design
directly enables these mechanisms by making the stable prefix structurally
explicit---the IMMUTABLE\_CONTEXT and CAPABILITY\_DECLARATION layers are
always positioned first and never change across invocations in a session,
making them natural cache targets.

\paragraph{Formal specifications for LLMs.}
Closest to ICS in spirit is TypeChat~\citep{typechat2023}, which uses type
schemas to constrain outputs, and work on schema-constrained
decoding~\citep{willard2023efficient,geng2023grammar}. These focus on
\emph{output} structure; ICS defines contracts for both input and output.
There is, to our knowledge, no prior work that defines a formal
\emph{multi-layer input specification} for LLM instructions with explicit
lifetime semantics, directive syntax, and a conformance model.

% ─────────────────────────────────────────────────────────────────────────────
\section{The ICS Specification}
\label{sec:spec}

An ICS-compliant instruction is composed of exactly five layers, appearing
in the canonical order defined below. Each layer is delimited by boundary
markers:

\begin{lstlisting}
###ICS:<LAYER_NAME>###
<layer content>
###END:<LAYER_NAME>###
\end{lstlisting}

Layers MUST appear in order. A parser encountering layers out of order MUST
reject the instruction as non-conformant. The five layers and their lifetime
semantics are as follows.

\subsection{IMMUTABLE\_CONTEXT}
\label{sec:ic}

IMMUTABLE\_CONTEXT declares long-lived facts that remain constant across all
tasks within a domain: system description, repository or schema layout,
architectural invariants, external dependency contracts. It MUST NOT include
task-specific instructions. It is the only layer eligible for aggressive
prompt caching, because it is stable across sessions by definition.

\textbf{Key rules.} Content MUST NOT be restated, redefined, or contradicted
in any subsequent layer. Any content that would be valid in IMMUTABLE\_CONTEXT
but is placed in SESSION\_STATE or TASK\_PAYLOAD is a conformance violation.

\subsection{CAPABILITY\_DECLARATION}
\label{sec:cd}

CAPABILITY\_DECLARATION defines what the model is permitted, forbidden, and
required to do. All permissions are expressed using three directives:
\texttt{ALLOW}, \texttt{DENY}, and \texttt{REQUIRE}. Undeclared actions are
DENIED by default.

\subsubsection{Directive grammar}
\label{sec:capgram}

\begin{lstlisting}[language=ICS]
directive  ::= KEYWORD WS+ action (WS+ qualifier)? (WS+ "IF" WS+ condition)?
KEYWORD    ::= "ALLOW" | "DENY" | "REQUIRE"
action     ::= WORD (WS+ WORD)*
qualifier  ::= QWORD WS+ target
QWORD      ::= "WITHIN" | "ON" | "WITH" | "UNLESS"
target     ::= WORD (WS+ WORD)*
condition  ::= WORD (WS+ WORD)*
WORD       ::= [A-Za-z0-9_./-]+
WS         ::= " " | "\t"
\end{lstlisting}

The grammar is intentionally liberal on scope expressions to support existing
codebase idioms without requiring rewrites at adoption time. A normative scope
grammar is deferred to ICS v0.2 as an optional conformance level. Example
directives:

\begin{lstlisting}[language=ICS]
ALLOW   file modification WITHIN src/orders/
DENY    modification of src/orders/api/
DENY    introduction of float arithmetic ON monetary values
REQUIRE type annotations ON all new functions
REQUIRE backward compatibility WITH gateway/v2/ API contracts
\end{lstlisting}

\subsection{SESSION\_STATE}
\label{sec:ss}

SESSION\_STATE captures temporary assumptions and decisions that accumulate
within a session but must not outlast it: conclusions reached in prior
invocations, temporary flags, confirmed facts about the current task scope.
Each entry SHOULD carry an identifier or timestamp.

A session ends when SESSION\_STATE contains only the sentinel value
\texttt{CLEAR}:

\begin{lstlisting}
###ICS:SESSION_STATE###
CLEAR
###END:SESSION_STATE###
\end{lstlisting}

A SESSION\_STATE layer containing \texttt{CLEAR} alongside other content MUST
be rejected as malformed. Declaring a new IMMUTABLE\_CONTEXT implicitly clears
any prior SESSION\_STATE.

\subsection{TASK\_PAYLOAD}
\label{sec:tp}

TASK\_PAYLOAD is the only layer that changes per invocation. It specifies the
precise action to be performed, without restating context from preceding
layers. A TASK\_PAYLOAD is executable without clarification if all required
inputs, constraints, and success conditions are fully defined across the
instruction layers.

\subsection{OUTPUT\_CONTRACT}
\label{sec:oc}

OUTPUT\_CONTRACT declares the required shape of the response. It is response-
scoped but MUST be declared before execution. All four fields are required:

\begin{center}
\begin{tabularx}{\linewidth}{lX}
\toprule
\textbf{Field} & \textbf{Meaning} \\
\midrule
\texttt{format}     & Output encoding (JSON, unified diff, prose, \ldots) \\
\texttt{schema}     & Structure definition (JSON Schema, field list, \ldots) \\
\texttt{variance}   & Explicitly permitted deviations; unlisted variance is forbidden \\
\texttt{on\_failure} & Required behavior when output cannot conform \\
\bottomrule
\end{tabularx}
\end{center}

Free-form prose is invalid unless \texttt{format: prose} is declared
explicitly. Outputs that do not satisfy the contract MUST NOT be partially
accepted; the \texttt{on\_failure} behavior MUST be applied.

\subsection{Conformance}
\label{sec:conform}

An instruction is \emph{ICS-compliant} if and only if all five layers are
present, delimited correctly, in canonical order, and satisfy the layer-
specific rules above. Partial compliance is not defined; the predicate is
binary.

A conformance checker MUST verify, in order: (1) all five boundary tags are
well-formed and present; (2) layers appear in canonical order; (3) SESSION\_STATE
containing \texttt{CLEAR} contains no other content; (4) no layer restates,
redefines, or contradicts a preceding layer; (5) CAPABILITY\_DECLARATION uses
only \texttt{ALLOW}/\texttt{DENY}/\texttt{REQUIRE} directives; (6) OUTPUT\_CONTRACT
contains all four required fields.

% ─────────────────────────────────────────────────────────────────────────────
\section{Worked Examples}
\label{sec:examples}

We provide fully conformant examples across five production-representative
domains to demonstrate that ICS is not specific to software development but
applies wherever an LLM is given repeatable, structured tasks. Abbreviated
examples appear below; complete examples are in the accompanying appendices.

\paragraph{Domain 1: Code refactoring.}
A development team drives automated modifications on an order management
service. IMMUTABLE\_CONTEXT declares the language, repository layout, and
monetary invariants; CAPABILITY\_DECLARATION locks the API layer against
modification; SESSION\_STATE accumulates confirmed findings about the
discount function; TASK\_PAYLOAD requests a targeted split of one function.

\paragraph{Domain 2: Log analysis.}
A recurring analytics pipeline runs latency breakdowns against API gateway
logs. CAPABILITY\_DECLARATION includes \texttt{DENY output of individual
caller\_id values} to enforce a privacy constraint that would otherwise be
re-stated (and potentially forgotten) in each TASK\_PAYLOAD.

\paragraph{Domain 3: Infrastructure change (Terraform).}
A platform team manages multi-account AWS infrastructure. CAPABILITY\_DECLARATION
uses absolute \texttt{DENY} directives to lock \texttt{infra/prod/} against
modification regardless of task phrasing, and \texttt{REQUIRE tagging} encodes
a cost-attribution convention that applies to every task in the session.

\paragraph{Domain 4: Database schema review.}
A review-only session assesses Alembic migration scripts.
CAPABILITY\_DECLARATION declares \texttt{DENY generation of migration file
content}---the model's role is reviewer, not author---and three
\texttt{REQUIRE} directives encode the team's known error classes. The
OUTPUT\_CONTRACT includes \texttt{verdict: "PASS" | "FAIL"} as an enumerated
field, making the pass/fail signal machine-checkable without inspecting the
violations array.

\paragraph{Domain 5: B2B payments platform (RAG + agent).}
A real-world example used in our empirical evaluation (\S\ref{sec:eval}).
IMMUTABLE\_CONTEXT covers 15 architectural invariants, 9 core data models,
and a full event-type registry. CAPABILITY\_DECLARATION includes
compliance-critical directives (\texttt{DENY logging of PII fields},
\texttt{DENY introduction of float arithmetic ON monetary values}).
The permanent context is 4{,}115 API-counted tokens---sufficient to trigger
the prompt-caching threshold on \texttt{claude-haiku-4-5-20251001}.

% ─────────────────────────────────────────────────────────────────────────────
\section{Evaluation}
\label{sec:eval}

We evaluate ICS along two axes: (1) \emph{structural token savings}---the
reduction in tokens sent per invocation attributable to lifetime separation
alone, independent of any pricing scheme; and (2) \emph{effective cost
savings}---the combined effect of structural savings and prompt-caching
pricing on real API responses.

\subsection{Structural savings: mathematical proof}
\label{sec:proof}

Let $T$ be the total token count of a full ICS instruction.
Let $P$ be the token count of the permanent layers
(IMMUTABLE\_CONTEXT $+$ CAPABILITY\_DECLARATION).
Let $V = T - P$ be the token count of the variable layers
(SESSION\_STATE $+$ TASK\_PAYLOAD $+$ OUTPUT\_CONTRACT).

The na\"{i}ve approach re-sends $T$ tokens on every invocation:
\[
  C_{\text{naive}}(N) = T \cdot N
\]

The ICS approach sends permanent layers once (or on session reset) and
variable layers on every invocation. Let $S \leq N$ be the number of
session-state changes. Then:
\[
  C_{\text{ics}}(N) = P \cdot 1 \;+\; V_S \cdot S \;+\; V_T \cdot N
\]
where $V_S$ and $V_T$ are the session and task/output token counts
respectively. Because $P > 0$ and $P < T$, we have
$C_{\text{ics}}(N) < C_{\text{naive}}(N)$ for all $N > 1$.

This is an identity---it holds for every ICS-compliant instruction regardless
of content or model. The savings rate converges to $P/T$ as $N \to \infty$.

\subsection{Token counting: method independence}
\label{sec:counting}

Table~\ref{tab:counting} shows token counts and savings at $N{=}10$
for two offline counting methods applied to the code-refactoring example
from Appendix A.

\begin{table}[h]
\centering
\caption{Token savings at $N{=}10$ across two counting methods.}
\label{tab:counting}
\begin{tabular}{llrr}
\toprule
\textbf{Example} & \textbf{Method} & \textbf{Na\"{i}ve tokens} & \textbf{Savings} \\
\midrule
Code refactoring & Approx ($\ell/4$) & 4,830 & 53.7\% \\
Code refactoring & Word-boundary BPE & 4,830 & 55.2\% \\
Log analysis     & Approx ($\ell/4$) & 5,520 & 55.2\% \\
Log analysis     & Word-boundary BPE & 5,520 & 53.1\% \\
\bottomrule
\end{tabular}
\end{table}

The savings percentage is stable at 53--55\% across both examples and both
methods, confirming that the structural saving is not an artifact of any
particular token-counting approximation.

\subsection{Savings as a function of invocation count}
\label{sec:scaling}

Table~\ref{tab:scaling} shows how savings grow as $N$ increases on the
code-refactoring example (word-boundary counting). Savings grow
monotonically, approaching $P/T \approx 63\%$ as $N \to \infty$.

\begin{table}[h]
\centering
\caption{Token savings as a function of invocation count $N$.}
\label{tab:scaling}
\begin{tabular}{rrrrr}
\toprule
$N$ & Na\"{i}ve tokens & ICS tokens & Saved & Savings \% \\
\midrule
1  & 483    & 483    & 0      & 0.0\%  \\
5  & 2,415  & 1,313  & 1,102  & 45.6\% \\
10 & 4,830  & 2,166  & 2,664  & 55.2\% \\
20 & 9,660  & 3,872  & 5,788  & 59.9\% \\
50 & 24,150 & 8,990  & 15,160 & 62.8\% \\
\bottomrule
\end{tabular}
\end{table}

\subsection{Live API measurement}
\label{sec:live}

We ran both the na\"{i}ve and ICS approaches against the Anthropic Messages
API using \texttt{claude-haiku-4-5-20251001} on the payments-platform example
(Domain~5). The ICS approach places the permanent layers in a single content
block marked \texttt{cache\_control: ephemeral}; the variable layers are in
separate unmarked blocks. The na\"{i}ve approach sends the full instruction
as a flat string with no cache markup.

We measured \texttt{input\_tokens}, \texttt{cache\_creation\_input\_tokens},
and \texttt{cache\_read\_input\_tokens} from the API response \texttt{usage}
field across 10 invocations.

\begin{table}[h]
\centering
\caption{Live API token usage at $N{=}10$. Permanent layer: 4,115 tokens.}
\label{tab:live}
\begin{tabular}{lrrrr}
\toprule
\textbf{Approach} & \textbf{Full-rate tokens} & \textbf{Cache-write} & \textbf{Cache-read} & \textbf{Cost (USD)} \\
\midrule
Na\"{i}ve & 46,040 & 0      & 0      & \$0.03811 \\
ICS       & 4,850  & 0      & 41,150 & \$0.00845 \\
\bottomrule
\end{tabular}
\end{table}

Applying Anthropic's caching rates (cache writes at $1.25\times$ standard,
cache reads at $0.10\times$ standard), the ICS approach achieves a
\textbf{77.8\%} cost reduction at $N{=}10$ (\$0.02966 saved per 10-invocation
session). The cache was warm from invocation 1 due to a prior warm-up call
during threshold testing; every subsequent ICS invocation read 4,115 tokens
from cache at the reduced rate.

\paragraph{Cache activation threshold.}
We empirically established that \texttt{claude-haiku-4-5-20251001} requires
$\geq{\sim}4{,}096$ tokens in the cached block to activate caching---higher
than the 1,024-token threshold documented for Claude~3 models. ICS's
requirement that IMMUTABLE\_CONTEXT be semantically complete encourages callers
to include full context in this layer, naturally pushing its size toward the
activation threshold for real-world workloads.

% ─────────────────────────────────────────────────────────────────────────────
\section{Discussion}
\label{sec:discussion}

\subsection{Why the output contract matters}
\label{sec:oc-discuss}

Requiring an OUTPUT\_CONTRACT before execution forces the caller to
articulate success conditions prior to invoking the model. This has two
practical effects. First, it makes failure states machine-detectable: a
caller checking for the \texttt{BLOCKED:} prefix or the
\texttt{\{"status": "error"\}} JSON envelope does not need to parse the
output to determine whether the task succeeded. Second, it shifts the retry
decision: if an output fails the contract check, the correct response is to
update the relevant layer, not to re-prompt with corrective prose. This
distinction---interface update versus retry---is the core behavioral
difference between ICS-compliant systems and ad hoc prompting.

\subsection{Adoption considerations}
\label{sec:adoption}

ICS is more verbose than informal prompting. A five-layer instruction for a
simple task may be longer than the equivalent prose prompt. This overhead is
intentional and appropriate for systems where instructions are authored once
and executed many times---the cost of verbosity is paid once, while the
benefits (caching, auditability, maintainability) compound across invocations.
For one-off or exploratory use, ICS is not the right tool.

The scope grammar in CAPABILITY\_DECLARATION is intentionally liberal in
v0.1. Two implementations may express the same constraint using different
syntax and both claim compliance. This trades machine-parseability for
adoption breadth: teams can migrate their existing constraint idioms without
rewriting. A normative scope grammar is planned for v0.2 as an optional
conformance level.

\subsection{Complementary to existing approaches}
\label{sec:complementary}

ICS is not a replacement for prompt optimization (DSPy), schema-constrained
decoding (Guidance, LMQL), or context management systems (MemGPT, RAG). It
operates at the layer \emph{below} these: it defines how an instruction is
structured and validated, not how it is generated, optimized, or executed.
A DSPy pipeline could emit ICS-compliant instructions; a RAG system could
place retrieved documents in TASK\_PAYLOAD while keeping system context in
IMMUTABLE\_CONTEXT.

\subsection{Limitations}
\label{sec:limits}

\textbf{No user study.} We demonstrate token and cost savings empirically but
do not report a user study on whether teams find ICS usable or whether it
reduces prompt-related bugs in practice. Adoption evidence beyond synthetic
examples remains future work.

\textbf{Single model family tested.} Live experiments use \texttt{claude-haiku-4-5-20251001}.
The structural savings proof is model-agnostic, but the observed
cost savings depend on a provider's caching pricing schedule. Experiments on
OpenAI and other providers with caching APIs are deferred to future work.

\textbf{Multi-modal context deferred.} ICS v0.1 treats instruction content as
text. Images, audio, and tool call results that should participate in lifetime
separation are not addressed. Multi-modal context lifetime is deferred to
v0.2.

\textbf{Session persistence is caller-managed.} ICS defines session semantics
(begin, accumulate, clear) but does not define storage mechanisms. Callers
must maintain SESSION\_STATE across invocations externally. ICS does not
provide a persistence protocol.

% ─────────────────────────────────────────────────────────────────────────────
\section{Conclusion}
\label{sec:conclusion}

We have presented ICS, a formal five-layer specification for LLM instructions
that treats prompts as interfaces rather than prose. ICS separates instruction
content by lifetime, declares capabilities using verifiable directive syntax,
and requires explicit output contracts before execution. These design choices
address the four failure modes that recur in production LLM systems---context
collapse, implicit constraints, ambiguous output expectations, and retry as
a first resort---at the structural level rather than through better prompting
practices.

Empirically, ICS's lifetime separation enables prompt-caching APIs to reduce
effective cost by 77.8\% at $N{=}10$ on a real-world B2B payments workload.
The underlying token savings (55.2\% structural, independent of pricing) are
proven to hold for any ICS-compliant instruction and any counting method.

The specification, validator, token analyzer, live test harness, and five
domain examples are available at
\url{https://github.com/[repository]}.

% ─────────────────────────────────────────────────────────────────────────────
\section*{Acknowledgements}
[Acknowledgements redacted for review.]

% ─────────────────────────────────────────────────────────────────────────────
\bibliographystyle{plainnat}
\bibliography{references}

\end{document}
